% Generated draft: requires mathematical and source review.
\documentclass[11pt]{article}
\usepackage[margin=25mm]{geometry}
\usepackage{fontspec}
\setmainfont{Latin Modern Roman}
\usepackage{amsmath,amssymb,mathtools,unicode-math}
\setmathfont{Latin Modern Math}
\usepackage{enumitem,needspace,xurl,hyperref}
\setlength{\parindent}{0pt}
\setlength{\parskip}{5pt}
\setlength{\emergencystretch}{3em}
\newcommand{\N}{\mathbb{N}}
\newcommand{\Z}{\mathbb{Z}}
\newcommand{\Q}{\mathbb{Q}}
\newcommand{\R}{\mathbb{R}}
\newcommand{\C}{\mathbb{C}}
\newcommand{\F}{\mathbb{F}}
\newcommand{\A}{\mathbb{A}}
\newcommand{\PP}{\mathbb{P}}
\newcommand{\E}{\mathbb{E}}
\newcommand{\eps}{\varepsilon}
\newcommand{\dd}{\,\mathrm{d}}
\begin{document}
\section*{UnsolvedMath ID: 11100027}
\textbf{Source problem code:} catalogue entry
\subsection{Definitions and mathematical statement}
\emph{Applications research problem (Mark Hovey's list).} Investigate Voevodsky's stable homotopy category of schemes from a homotopy theorist's point of view, with particular attention to Bousfield localization. The source explicitly calls this a huge, unstructured program rather than a specified conjecture.

For a chosen base scheme \(B\) in an appropriate setup, a modern conventional model is the motivic stable homotopy category \(\mathcal{SH}(B)\), built from pointed motivic spaces by imposing the chosen local descent and \(\A^1\)-homotopy equivalences and stabilizing with respect to the motivic sphere. The excerpt does not specify a base, topology, stabilization convention, or foundational hypotheses, so this is explanatory notation rather than an unprovided fixed setting.

A Bousfield localization should be defined by the selected equivalences or acyclic objects, with a localization functor \(L\) and unit making \(LX\) universal among local objects receiving a map from \(X\). A precise instance of the research program must choose which localization, characterize its local objects and effect on the intended motivic invariants, and state any compatibility with the monoidal or stable structure. The source chooses no spectrum, map class, prime, or specific theorem, so no particular chromatic, slice, or other construction is identified as its unique target.

Hovey says he believed localization had not yet been exploited and notes the algebraic-geometric background as a practical obstacle. These are historical personal assessments, not current mathematical assertions. The summary calls \(\mathcal{SH}\) and several localization programs mature and the broader program partly realized but still open; it supplies no theorem or completion criterion. Slice constructions, chromatic localizations, and an undefined ``synthetic theory'' are not silently treated as identical notions.
\subsection{Short English statement}
Develop homotopy-theoretic approaches to Voevodsky's stable category of schemes, especially Bousfield localizations. The source deliberately gives a broad program, without a base scheme or chosen localization, and its no-prior-use remark is historical.
\subsection{Sources}
\begin{itemize}
\item Local UnsolvedMath registry, permanent ID 11100027, source code catalogue entry; proposer Mark Hovey, ``Applications,'' item 3, explicitly referring to Voevodsky's stable homotopy category.
\item \emph{Hovey's Algebraic Topology Problem List}, surviving mirror: \url{https://www-users.cse.umn.edu/~tlawson/hovey/}. Recorded access 2026-07-29; extraction html-ordered-list. The mirror has 65 top-level items across nine category pages.
\item No precise theorem or full external citation is supplied for the category or a proposed localization. The summary names chromatic/motivic, \(K(1)\)-local, and slice directions as modern progress, without specifying their settings.
\item Supplied status partially\_solved and PARTIAL-PROGRESS. Its summary says largely realized with a broad still-open remainder, an editorial assessment without a precise completion criterion. Current primary-literature status and redistribution terms remain NEEDS\_REVIEW. The registry assigns L4 whereas background records default L3.
\end{itemize}
\end{document}
